\chapter{\textit{Stationary Distributions in IDTMCs}}
\label{ch:BSCCs_in_IDTMCs}
% What is a stationary distribution in the context of Markov Chains?
% --> If we view vector \pi_i as the vector of prob. distribution in step i, then stationary distribtion, is when this vector remains unchanged despite mutiplication with transition matrix

% Introduce concept of stationary distribution
In this chapter, we examine how stationary distributions in IDTMCs can be calculated using EVTs. We start by formally introducing stationary distributions for DTMCs and IDTMCs. We then continue by lifting results from \cite{modeling_and_analysis_stochastic_systems,Computing_EVTs_Stationary_Distr_in_MCs_Mertens} and extending them to IDTMCs. To this end, we fix an IDTMC $\IMC = (S^\IMC,\PFlower,\PFupper,\Initlower,\Initupper)$ and a DTMC $\MC = (S^\MC,\PF,\iota_{init})$ with $\MC \models_U \IMC$.
\todo{fix DTMC and IDTMC}

We define \textit{stationary distributions}, also called long-run distributions, as done in \cite{Bibel}.
\begin{definition}[Stationary Distributions in DTMCs]
    Given a finite DTMC $\MC = (S,\PF,\iota_{init})$ and $s\in S$, the \textit{stationary distribution} $\theta^\MC(s)$ is given by
    \[ \theta^\MC(s) = \lim_{n\rightarrow \infty} \frac{1}{n} \cdot \sum_{i=1}^{n}Pr(\lozenge^{=k}s)\ .\]
    % The function $t \mapsto \theta(s,t)$ is the \textit{stationary distribution} for starting state $s$. The stationary distribution for $\MC$ is given by
    % \[ \Theta^\MC(t) = \sum_{s\in S} \iota_{init}(s) \cdot \theta(s,t)\ .\]
\end{definition}
In other words, the stationary distribution $\theta^\MC(s)$ describes the probability that $\MC$ is in state $s$ after $k$ steps, as $k$ tends towards infinity. 
%Whenever the DTMC $\MC$ is evident from the context, we omit the superscript $\MC$ from $\theta^\MC$.

Similarly, we now define stationary distributions for IDTMCs. First, recall that the DTMCs $\MC_{\eta} \models_I \IMC$ that are induced by some nature $\eta$ can be infinite, i.e., $|\MC_{\eta}| = \infty$, see \Cref{def:DTMC_induced_by_nature}. Intuitively, as every state $s\in S^{\MC_{\eta}}$ in these DTMCs is only visited once, it follows that $\theta(s) = 0$. 
Thus, we restrict this analysis of stationary distributions in this thesis to the UMC semantic.

\begin{definition}[Stationary Distributions in IDTMCs]
    For a given IDTMC $\IMC$ and state $s\in S^\IMC$, we define 
    \[\Theta(s) = \{\theta^\MC(s) \mid \MC \models_U \IMC\}\]
    as the set of stationary distributions of $s$ on $\IMC$. 
    \\Furthermore, we define $\Thetalower^\IMC(s)$ and $\Thetaupper^\IMC(s)$ such that 
    \[\Thetalower^\IMC(s) = \operatorname{inf}_{\MC\models_U\IMC}\theta^\MC(s) \qquad \text{and} \qquad \Thetaupper^\IMC(s) = \operatorname{sup}_{\MC\models_U\IMC}\theta^\MC(s)\ .\]  
    %\[\]
\end{definition}
In other words, $\Thetalower^\IMC(s)$ and $\Thetaupper^\IMC(s)$ describe the minimum and maximum stationary distribution of $s$ among the DTMCs $\MC \models_U \IMC$ that satisfy $\IMC$ under the UMC semantic.

Moreover, BSCCs are defined as follows for IDTMCs.
\begin{definition}[BSCCs for IDTMCs]
    Let $\IMC$ be an IDTMC, then $B^\IMC \subseteq S^\IMC$ is a BSCC if there exists at least one DTMC $\MC \models_U \IMC$ such that $B^\IMC$ is a BSCC in the DTMC $\MC$, see \Cref{def:(B)SCC}.
\end{definition}
Whenever the DTMC $\MC$ is evident from the context, we omit the superscript $\IMC$ from $B^\IMC$.
% \todo{stationary distribution gar nicht möglich für IMDPs?; weil induzierte DTMCs infinite sein können; oder gibt es dann einfach keine stationary distribution in den DTMCs}

In contrast to EVTs, we are interested in the recurrent and not the transient states when considering the stationary distribution of a DTMC or IDTMC. Recall that the recurrent states of a DTMC $\MC$ are precisely the states contained in the BSCCs of $\MC$, see \Cref{def:transient_and_recurrent_states}. Thus, by collapsing each BSCC of $\MC$ into a single state, the stationary distribution $\theta^\MC$ of the recurrent states of a DTMC $\MC$ is characterized by 
\begin{equation}\label{eq:Stat_Distr_via_BSCC_for_DTMC} 
    \theta^\MC(s) = Pr^\MC(\lozenge B) \cdot \theta^{\MC|_B}(s)\ ,
\end{equation} 
for all $s\in B$, as proven in \cite{modeling_and_analysis_stochastic_systems}, where $\theta^{\MC|_B}$ denotes the stationary distribution of $\MC$ restricted to $B$. Moreover, in \cite{Computing_EVTs_Stationary_Distr_in_MCs_Mertens}, it is shown that the stationary distribution within a BSCC can be computed through its EVTs. To this end, they construct a DTMC $\mathcal{B}\!\Rsh^{\hat{v}}$ for every BSCC $B$ with any arbitrary state $v\in S^{\MC|_B}$ in which all incoming transitions of $v$ are redirected to the fresh absorbing state $\hat{v}$. As $v$ is arbitrary, we fix $v$ and $\hat{v}$ for the remainder of this chapter.
\todo{vllt reicht das nicht: muss nochmal genauer argumentieren, weshalb auch für IDTMC $v$ arbiträr}

\emph{Construction.} Formally, $\mathcal{B}\!\Rsh^{\hat{v}} = (S^{B} \cup \{\hat{v}\},\PF^{\mathcal{B}\Rsh^{\hat{v}}},\iota_{init}^{\mathcal{B}\Rsh^{\hat{v}}})$, where $S^B$ denotes the states of $\MC$ that are contained in $B$. Moreover, $\PF^{\mathcal{B}\Rsh^{\hat{v}}}(\hat{v},\hat{v}) = 1$, $\PF^{\mathcal{B}\Rsh^{\hat{v}}}(s,\hat{v}) = \PF^\MC(s,v)$, and $\PF^{\mathcal{B}\Rsh^{\hat{v}}}(s,v) = 0$ for all $s\in S^{\MC|_B}$. Lastly, $\PF^{\mathcal{B}\Rsh^{\hat{v}}}(s,t) = \PF^\MC(s,t)$ for all $s\in S$, $t\in S\backslash\{v\}$, and $\iota_{init}^{\mathcal{B}\Rsh^{\hat{v}}}(v) = 1$. 

The stationary distribution within a DTMC is then given by
\begin{equation}\label{eq:Stat_Distr_of_BSCC}
    \theta^{\MC|_B}(s) = \frac{\Exp^{\mathcal{B}\Rsh^{\hat{v}}}[vt_s]}{\sum_{t\in S}\Exp^{\mathcal{B}\Rsh^{\hat{v}}}[vt_t]}\ ,
\end{equation}
as proven in \cite{Computing_EVTs_Stationary_Distr_in_MCs_Mertens}.

Similarly to the DTMC $\mathcal{B}\!\Rsh^{\hat{v}}$, we now define the IDTMC $\mathcal{I}\!\Rsh^{\hat{v}}$ as follows.

\emph{Construction.} $\mathcal{I}\!\Rsh^{\hat{v}} = (S^{B} \cup \{\hat{v}\},\PFlower^{\mathcal{B}\Rsh^{\hat{v}}},\PFupper^{\mathcal{B}\Rsh^{\hat{v}}},\Initlower^{\mathcal{B}\Rsh^{\hat{v}}},\Initupper^{\mathcal{B}\Rsh^{\hat{v}}})$, where $S^B$ denotes the states that are contained in $B$. Moreover, $\PFupper^{\mathcal{B}\Rsh^{\hat{v}}}(\hat{v},\hat{v}) = \PFlower^{\mathcal{B}\Rsh^{\hat{v}}}(\hat{v},\hat{v}) = 1$, $\PFupper^{\mathcal{B}\Rsh^{\hat{v}}}(s,\hat{v}) = \PFupper^\IMC(s,v)$, $\PFlower^{\mathcal{B}\Rsh^{\hat{v}}}(s,\hat{v}) = \PFlower^\IMC(s,v)$, and $\PFupper^{\mathcal{B}\Rsh^{\hat{v}}}(s,v) = \PFlower^{\mathcal{B}\Rsh^{\hat{v}}}(s,v) = 0$ for all $s\in S^{\MC|_B}$. Lastly, $\PF^{\mathcal{B}\Rsh^{\hat{v}}}(s,t) = \PF(s,t)$ for all $s\in S$, $t\in S\backslash\{v\}$, and $\iota_{init}^{\mathcal{B}\Rsh^{\hat{v}}}(v) = 1$. 

Intuitively, as $\MC \models_U \IMC$ it follows, by construction, that $\mathcal{B}\!\Rsh^{\hat{v}} \models_U \mathcal{I}\!\Rsh^{\hat{v}}$.

% Es gibt IDTMCs nicht unbedingte feste BSCCs, i.e. s kann in verschiedenen BSCCs sein.

\begin{assumption}\label{assmp:graph_preserving}
    In the rest of this chapter, we assume that all IDTMCs $\IMC$ are graph preserving, meaning that there exists no lower transition probability function such that $\PFlower^\IMC(s,t) = 0$ for any $s,t\in S^\IMC$. This implies that there exists no $\MC \models_U \IMC$ such that some $s\in S^\IMC$ is unreachable in $\MC$.
\end{assumption}

Given this assumption, we can prove the following theorem.
\begin{theorem}
    Given an IDTMC $\IMC$, a state $s\in S_{re}^\IMC$, and any pair of BSCCs $B^{\MC},B'^{\MC'} \subseteq S^\IMC$ with $B \neq B'$, it holds that $s\in B$ implies $s\notin B'$.
\end{theorem}
In other words, every state $S_{re}^\IMC$ corresponds to precisely one BSCC in $\IMC$.

\begin{proof}
    Assume that for some IDTMC $\IMC$ there exists a state $s$, two BSCCs $B^{\MC}, B'^{\MC'}$ with $B^\neq B'^{\MC'}$, and $\MC, \MC'\models \IMC$, such that $s\in B^{\MC}$ and $s\in B'^{\MC'}$. As $B^{\MC}\neq B'^{\MC'}$, there exists at least one state $t\in S^\IMC$ that is only contained in either $B^{\MC}$ or $B'^{\MC'}$, i.e., $t\in (B\backslash B') \cup (B'\backslash B)$. 
    
    Without loss of generality, let us assume $t\in B^\MC$. Hence, by \Cref{def:(B)SCC}, there exists a path from $s$ to $t$ in $\MC$. However, as noted, the graph relations for all satisfying DTMCs of graph preserving IDTMCs are identical. Thus, there must also exist a path from $s$ to $t$ in $\MC'$. However, $t\notin B'^{\MC'}$, which contradicts \Cref{def:(B)SCC}.
\end{proof}

\todo{BSCC mit rücksicht auf IMC definieren, also dass BSSC von DTMC, die IDTMC satisfied}

Thus, we can conclude for graph preserving IDTMCs that any state $s\in S^\IMC$ belongs only to a singular BSCC $B$, i.e., the BSCCs for all satisfying DTMCs are identical. This allows us to optimize the reachability of $B$ ($Pr^\MC(\lozenge B)$) and the stationary distribution restricted to $B$ ($\theta^{\MC|_B}$) of \Cref{eq:Stat_Distr_via_BSCC_for_DTMC} independently as shown in the following.

Therefore, we can intuitively calculate the stationary distributions for a state $s\in \IMC$ by solving the following NLP:

\[
   \begin{aligned}
        \operatorname{opt} \quad & \frac{\Exp^{\mathcal{B}\Rsh^{\hat{v}}}[vt_s]}{\sum_{t\in S^{\mathcal{B}\Rsh^{\hat{v}}}}\Exp^{\mathcal{B}\Rsh^{\hat{v}}}[vt_t]} & & \operatorname{opt}\in\{\max,\min\}\\
        \text{s.t.} \quad & \mathcal{B}\Rsh^{\hat{v}} \models_U \mathcal{I}\Rsh^{\hat{v}}\ .
    \end{aligned}
\]

\begin{theorem}\label{th:max_min_stat_distribution}
    Given an IDTMC $\IMC$, a state $s$, the corresponding BSCC $B^\IMC$ with $s\in B^\IMC$, and the corresponding $\mathcal{I}\Rsh^{\hat{v}}$, let $\overline{\theta}^{\IMC|_B}$ and $\underline{\theta}^{\IMC|_B}$ be the the solutions of the NLP for $\mathcal{I}\Rsh^{\hat{v}}$ when maximizing and minimizing, respectively.
    We can then calculate the minimum and maximum stationary distribution using
    \[\Thetalower^\IMC(s) = \Prlower^\IMC(\lozenge B) \cdot \underline{\theta}^{\IMC|_B}\]
    and 
    \[\Thetaupper^\IMC(s) = \Prupper^\IMC(\lozenge B) \cdot \overline{\theta}^{\IMC|_B}\]
    similar to \Cref{eq:Stat_Distr_via_BSCC_for_DTMC}.
\end{theorem}

We provide a short proof by contradiction of this theorem
\begin{proof}
    Let us assume that there exists a DTMC $\MC \models_U \IMC$ such that $\Thetaupper^\IMC(s) < \theta^\MC(s)$. Using \Cref{th:max_min_stat_distribution}, \Cref{eq:Stat_Distr_via_BSCC_for_DTMC}, and \Cref{eq:Stat_Distr_of_BSCC}, we can thus conclude that either
    \[Pr^\MC(\lozenge B) > \Prupper^{\IMC}(\lozenge B) \qquad \text{or} \qquad \frac{\Exp^{\mathcal{B}\Rsh^{\hat{v}}}[vt_s]}{\sum_{t\in S}\Exp^{\mathcal{B}\Rsh^{\hat{v}}}[vt_t]} > \overline{\theta}^{\IMC|_B}\ .\]
    However, as $\MC \models_U \IMC$ and $\mathcal{B}\!\Rsh^{\hat{v}} \models_U \mathcal{I}\!\Rsh^{\hat{v}}$, by the definition of $\Prupper^\IMC$ and the optimality of $\overline{\theta}^{\IMC|_B}$ as a maximal solution of the NLP, this leads to a contradiciton.
    Furthermore, as $s$ belongs to a singular BSCC $B$, every $\mathcal{B}\!\Rsh^{\hat{v}} \models_U \mathcal{I}\!\Rsh^{\hat{v}}$ is necessarily constructed from $B$. Thus, there exists a DTMC $\MC \models_U \IMC$, such that $Pr^\MC(\lozenge B) = \Prupper^\IMC(\lozenge B)$ and $\theta^{\MC|_B} = \overline{\theta}^{\IMC|_B}$.
    The proof for $\Thetalower^\IMC(s)$ follows analogously by inverting the inequalities.
\end{proof}

% Show: when $\Exp[vt_s]$ minimal then sum of $\Exp[vt_t]$ is maximal (and other way around). Bzw die summe aller evts ist für alle satisfying DTMCs gleich (--> nicht wahr; vllt auch ein interessantes Ergebnis??)

% This directly concludes the optimality of the BSCC restricted stat. distribution.



% (Trivially then uppper bound for stationary distribution given by uEVT and lower bound given by $\ell$EVT)

% \todo{Konstruktion rausnehmen??}